\documentclass[11pt,letterpaper]{article}
\usepackage[margin=0.85in]{geometry}
\usepackage{lmodern}
\usepackage{hyperref}
\hypersetup{colorlinks=false, urlcolor=black, pdfborder=0 0 0}
\setlength{\parindent}{0pt}
\setlength{\parskip}{0.45em}
\raggedright
\pagestyle{empty}

\begin{document}

\noindent
{\bfseries Ke-Cheng Hsu (Clio)}\\
Ph.D.\ Candidate, Department of Economics, HKUST\\
(852) 6840\,9753 \quad $|$ \quad \href{mailto:khsu@connect.ust.hk}{khsu@connect.ust.hk}\\
\href{https://cliohsu1997.github.io/}{cliohsu1997.github.io} \quad $|$ \quad \href{https://github.com/cliohsu1997}{github.com/cliohsu1997}

\vspace{0.5em}
\noindent
October 22, 2026

\vspace{0.35em}
\noindent
Hiring Team\\
Morgan Stanley --- Fixed Income Division, Taiwan Institutional Sales, Hong Kong\\
Re: Analyst, Taiwan Institutional Sales (Fixed Income Division; JR044975)

\vspace{0.2em}
Dear Hiring Team,

I am writing to apply for the Analyst role on the Taiwan Institutional Sales team in Hong Kong. I grew up in Taiwan and came to Hong Kong for university; I am a native Taiwanese speaker, professionally fluent in English, and a Hong Kong permanent resident. I am now a Ph.D.\ candidate in Economics at HKUST (expected June 2027; I can finish earlier if preferred). A HK-based desk covering Taiwan institutional clients across FX, credit, and rates is exactly the setting where I want to put Mandarin--English communication and careful market analysis to work.

My interest in finance is longstanding. My first research project, at National Taiwan University, modeled a bank crisis on an interbank lending network as a negotiation between the regulator and the banks' creditors, and simulated how network structure affects the credibility of a no-bailout threat. My undergraduate minor was actuarial mathematics, and my Ph.D.\ coursework covers merger simulation, market entry and exit counterfactuals, and firm-conduct identification. I now want to apply that training to preparing market updates and trade ideas, supporting client discussions, and deepening coverage of Taiwan institutional investors.

The posting asks for clear analysis, client-facing communication, and comfort in a fast-paced markets environment. That matches how I already work. In joint work with Kei Kawai and Kohei Kawaguchi, who consult at Compass Lexecon and UTokyo Economic Consulting, I cleaned point-of-sale records from more than 1{,}000 stores into an analyzable dataset, estimated demand and costs under tight deadlines, and delivered figures, tables, and short written summaries for decision-makers. On Taobao apparel e-commerce I automated data extraction with a scraper, an API, and AI-based OCR, then estimated a structural model of purchase and return decisions in a tested Python, R, and C++ pipeline. I use AI tools experimentally and responsibly: every pipeline carries tests, so any number traces back to the code that produced it. Teaching ratings of 4.72/5 (TA) and 4.53/5 (lecturer), and conference presentations, have taught me to explain a theme clearly and defend a number under questioning.

I am keen to learn FX, credit, and rates coverage on the desk, build relationships with Taiwan asset managers, insurers, and banks, and partner with trading, research, and syndicate colleagues. I learn new product and market settings quickly, I manage multiple priorities under deadlines, and I would welcome the chance to contribute to your Taiwan Sales team.

I have attached my resume. Thank you for your consideration.

\vspace{0.4em}
Sincerely,\\[0.8em]
Ke-Cheng Hsu (Clio)

\end{document}
